\begin{exercise}[Kernel restrictions]\label{ex:vi15-01}
Prove directly that $U\mapsto\ker\varphi_U$ is a presheaf and that every restriction map is the restriction inherited from $\mathcal F$.
\end{exercise}

\begin{exercise}[Kernel sheaf]\label{ex:vi15-02}
Prove the sheaf axiom for $\ker\varphi$ without using the theorem in the text.
\end{exercise}

\begin{exercise}[Kernel stalk]\label{ex:vi15-03}
Construct the canonical isomorphism $(\ker\varphi)_x\cong\ker\varphi_x$.
\end{exercise}

\begin{exercise}[Image presheaf]\label{ex:vi15-04}
Show that $U\mapsto\im\varphi_U$ is a subpresheaf of $\mathcal G$.
\end{exercise}

\begin{exercise}[Local-image restrictions]\label{ex:vi15-05}
Show that a section locally in the image remains locally in the image after restriction.
\end{exercise}

\begin{exercise}[Local-image gluing]\label{ex:vi15-06}
Prove that locally-in-the-image sections satisfy the sheaf gluing axiom.
\end{exercise}

\begin{exercise}[Image stalk]\label{ex:vi15-07}
Prove $(\im\varphi)_x=\im\varphi_x$ directly from germs.
\end{exercise}

\begin{exercise}[Quotient presheaf]\label{ex:vi15-08}
Verify that $U\mapsto\mathcal F(U)/\mathcal F'(U)$ is a presheaf.
\end{exercise}

\begin{exercise}[Quotient stalk]\label{ex:vi15-09}
Write down explicitly the map from the stalk of the quotient presheaf to $\mathcal F_x/\mathcal F'_x$ and prove it is an isomorphism.
\end{exercise}

\begin{exercise}[Cokernel stalk]\label{ex:vi15-10}
Prove $(\coker\varphi)_x\cong\coker\varphi_x$.
\end{exercise}

\begin{exercise}[Quotient kernel]\label{ex:vi15-11}
For $q:\mathcal F\to\mathcal F/\mathcal F'$, prove $\ker q=\mathcal F'$.
\end{exercise}

\begin{exercise}[Cokernel as quotient]\label{ex:vi15-12}
Prove $\coker\varphi\cong\mathcal G/\im\varphi$.
\end{exercise}

\begin{exercise}[Zero and identity]\label{ex:vi15-13}
Compute kernel, image, coimage, and cokernel for the zero morphism and the identity morphism.
\end{exercise}

\begin{exercise}[Multiplication by $n$]\label{ex:vi15-14}
Compute these four constructions for $n:\underline{\mathbb Z}\to\underline{\mathbb Z}$.
\end{exercise}

\begin{exercise}[Skyscraper map]\label{ex:vi15-15}
For a homomorphism $A\to B$, compute the kernel, image, and cokernel of the induced skyscraper-sheaf morphism at a closed point.
\end{exercise}

\begin{exercise}[Direct sums]\label{ex:vi15-16}
For the projection $\mathcal F\oplus\mathcal G\to\mathcal G$, identify kernel and cokernel.
\end{exercise}

\begin{exercise}[Diagonal]\label{ex:vi15-17}
Show that the cokernel of $\Delta:\mathcal F\to\mathcal F\oplus\mathcal F$ is isomorphic to $\mathcal F$.
\end{exercise}

\begin{exercise}[Restriction to an open]\label{ex:vi15-18}
Prove that image sheaves and quotient sheaves commute with restriction to an open subset.
\end{exercise}

\begin{exercise}[Local surjectivity]\label{ex:vi15-19}
Show that $\im\varphi=\mathcal G$ iff every $g\in\mathcal G(U)$ is locally liftable through $\varphi$.
\end{exercise}

\begin{exercise}[Stalkwise surjectivity]\label{ex:vi15-20}
Show that $\im\varphi=\mathcal G$ iff every stalk map $\varphi_x$ is surjective.
\end{exercise}

\begin{exercise}[Coimage]\label{ex:vi15-21}
Construct the canonical map $\mathcal F/\ker\varphi\to\im\varphi$.
\end{exercise}

\begin{exercise}[First isomorphism theorem]\label{ex:vi15-22}
Prove the map in the previous exercise is an isomorphism by taking stalks.
\end{exercise}

\begin{exercise}[Continuous exponential]\label{ex:vi15-23}
Show that $f\mapsto e^{2\pi i f}$ from real-valued to circle-valued continuous-function sheaves is locally surjective.
\end{exercise}

\begin{exercise}[No global argument]\label{ex:vi15-24}
Show that the identity $S^1\to S^1$ has no continuous lift $S^1\to\mathbb R$ through $t\mapsto e^{2\pi i t}$.
\end{exercise}

\begin{exercise}[Derivative and local primitives]\label{ex:vi15-25}
Show that $d:\mathcal C^\infty\to\Omega^1$ on a smooth one-manifold has image sheaf all of $\Omega^1$.
\end{exercise}

\begin{exercise}[Ideal-sheaf quotient]\label{ex:vi15-26}
If $\mathcal I\subseteq\mathcal O$ is a sheaf of ideals, prove that $\mathcal O/\mathcal I$ is naturally a sheaf of rings and identify its stalks.
\end{exercise}
